\documentclass[a4paper,12pt]{extarticle}
% Быстрый черновик при таймауте Overleaf: добавьте опцию draft к documentclass (плейсхолдеры вместо фигур).
\usepackage{geometry}
\usepackage[T1]{fontenc}
\usepackage[utf8]{inputenc}
\usepackage[english,russian]{babel}
\usepackage{amsmath}
\usepackage{amssymb}
\usepackage{setspace}
\usepackage{graphicx}
\usepackage{indentfirst}
\usepackage[
  backend=biber,
  style=numeric,
  maxbibnames=99
]{biblatex}
\addbibresource{refs.bib}
\usepackage{mathtools}
\usepackage{booktabs}
\usepackage{multirow}
\usepackage[flushleft]{threeparttable}
\usepackage{tablefootnote}

\usepackage{chngcntr} % нумерация графиков и таблиц по секциям
\counterwithin{table}{section}
\counterwithin{figure}{section}

\graphicspath{{graphics/}}%путь к рисункам

\usepackage[colorlinks,citecolor=blue,linkcolor=blue,bookmarks=false,hypertexnames=true,urlcolor=blue]{hyperref}

\makeatletter
% \renewcommand{\@biblabel}[1]{#1.} % Заменяем библиографию с квадратных скобок на точку:
\makeatother

\geometry{left=2.5cm,right=1.0cm,top=2.0cm,bottom=2.0cm}
\setlength{\parindent}{1.25cm}
\renewcommand{\baselinestretch}{1.5} % междустрочный интервал


\newcommand{\bibref}[3]{\hyperlink{#1}{#2 (#3)}} % biblabel, authors, year
\addto\captionsrussian{\def\refname{Список литературы (или источников)}} 

\renewcommand{\theenumi}{\arabic{enumi}}% Меняем везде перечисления на цифра.цифра
\renewcommand{\labelenumi}{\arabic{enumi}}% Меняем везде перечисления на цифра.цифра
\renewcommand{\theenumii}{.\arabic{enumii}}% Меняем везде перечисления на цифра.цифра
\renewcommand{\labelenumii}{\arabic{enumi}.\arabic{enumii}.}% Меняем везде перечисления на цифра.цифра
\renewcommand{\theenumiii}{.\arabic{enumiii}}% Меняем везде перечисления на цифра.цифра
\renewcommand{\labelenumiii}{\arabic{enumi}.\arabic{enumii}.\arabic{enumiii}.}% Меняем везде перечисления на цифра.цифра

\begin{document}
\input{title_vkr}% это титульный лист - выберите подходящий вам из имеющихся в проекте вариантов (kr - курсовая работа у 3 курса, vkr - выпускная квалификационная работа у 4 курса)
\newpage
\setcounter{page}{2}

{
	\hypersetup{linkcolor=black}
	\tableofcontents
}

\newpage

\newpage
\section*{Аннотация}   % this is how to use russian

\addcontentsline{toc}{section}{Аннотация}
В данной работе рассматривается задача сжатия контекста в текстовой модальности с минимальной потерей качества.

В качестве решения интересно исследовать подход на основе визуальных энкодеров, поскольку такие энкодеры уже обучены на выделение информации из визуального представления и хочется перенести эти знания в текстовую модальность. Архитектура такой модели похожа на современные VLM модели - она состоит из проекторов, которые приводят векторы скрытых состояний в размерность, с которой работает визуальный энкодер и другой, который проецирует в размерность текстового декодера. Соответственно второй компонент такой архитектуры - визуальный энкодер, который преобразует полученные скрытые представления так, чтобы текстовый декодер с их помощью мог восстановить текст. Также в архитектуру входит небольшой энкодер, позволяющий представить исходный текст с помощью некоторого количества эмбеддинг векторов. И последний компонент такой архитектуры - текстовый декодер, который по промпту и выходу энкодера после проектора, генерирует текст.  

В качестве данных для обучения и валидации используются датасеты, состоящие из промптов (задач) и длинных по количеству токенов ответов к ним.

В результате ожидаю получить модель, которая может сжимать приходящий текст и довольно точно восстанавливать сжатый текст. Далее обе части, отвечающие за сжатие и разжатие текста можно независимо использовать в различных задачах.

\section*{Ключевые слова}
Глубинное обучение, сжатие контекста, визуальные энкодеры, трансформеры, декодеры, LLM.
\pagebreak

\section{Введение} 

Большие языковые модели (LLM) демонстрируют выдающиеся результаты в самых разнообразных задачах обработки естественного языка: от ответов на вопросы и генерации текста до сложного многошагового рассуждения. Одним из ключевых факторов этого успеха является способность моделей принимать на вход длинный контекст, включающий инструкции, примеры, документы, результаты вызова инструментов (tool calling). Развитие таких технологий, как Chain-of-Thought, In-Context Learning и RAG, растущая популярность агентских систем, приводит к тому, что длина контекста неуклонно возрастает, нередко превышая десятки тысяч токенов, что влечёт за собой рост задержки инференса, превышение лимитов контекстного окна и существенное увеличение вычислительных затрат.

Благодаря наличию такой задачи активно развивается направление сжатия контекста. Существуют два основных подхода решения такой задачи - методы обучения непрерывных представлений контекста в пространстве эмбеддингов и инженерные методы, основанные на более эффективном сжатии KV-кешей, удалении малоинформативных токенов и максимальной утилизации Context Parallelism. К первому варианту, например, можно отнести In-context Autoencoder (ICAE)~\cite{icae} - модель, состоящую из текстового LLM-энкодера с обучаемым LoRA-адаптером для кодирования длинного контекста в небольшое число слотов памяти и замороженного LLM-декодера. Ко второму направлению относятся ClusterKV~\cite{cluster_kv}, где токены KV-кэша восстанавливаются на уровне семантических кластеров, и методы prompt compression, такие как Selective Attention Compression~\cite{sac}, в которых явный контекст сокращается за счёт выбора наиболее информативных токенов.

Вместе с тем, подходы, основанные на обучении эмбеддингов, как правило, используют в качестве энкодера тяжелые языковые модели, что сопряжено с высокими вычислительными затратами и высокими затратами памяти, как например, в вышеупомянутом методе ICAE.

Предметом исследования данной работы является применение визуальных энкодеров для построения компактных представлений текстового контекста. Ключевая гипотеза состоит в том, что визуальные энкодеры, обученные извлекать плотную информацию из изображения и размещать ее в ограниченном наборе векторов, обладают индуктивным смещением, которое может быть полезно в задаче сжатия текстовой информации, поскольку задача построения компактного представления большого объема данных едина вне зависимости от модальности. Архитектурно предлагаемый подход заимствует идею современных мультимодальных (VLM) моделей и включает текстовый энкодер, два модуля-проектора, визуальный энкодер и языковой декодер.

Целью работы является исследование применимости визуальных энкодеров для задачи сжатия текстового контекста и построить модель, способную сжимать текст в компактное латентное представление с последующим точным восстановлением исходного содержания декодером. Далее будет исследоваться возможность использования такого сжатого представления для применения в сложных задачах, требующих глубоких рассуждений, чтобы еще больше сжать количество используемых токенов.

Научная новизна работы состоит в постановке и исследовании задачи переноса знаний визуальных энкодеров в текстовую модальность для целей сжатия контекста.

Практическая значимость результатов состоит в возможности независимого применения обученных компонентов: энкодера - для эффективного сжатия длинного текста в RAG системах или при агентском взаимодействии и вызове инструментов, декодера - для корректного восстановления и верификации сжатой информации. Это дает возможность снизить вычислительные затраты при работе с длинными контекстами в существующих LLM-решениях. 


% Ссылки на статьи оформляются с помощью пакета \texttt{biblatex}, например~\cite{chirkova18}. В описании статье в bib файле нужно обязательно указывать место публикации работы (журнал или конференцию) и год. Обратите внимание, что для описания статей из разных источников в списке литературы используются разные команды в bib файле: статья из журнала~\cite{ctan}, статья с конференции~\cite{chirkova18}, книга~\cite{knuth-acp}, глава книги~\cite{knuth-fa}. Если статься еще не опубликована нигде, а только выложена на arXiv, то на нее тоже можно сослаться~\cite{chirkova18_arxiv}, но предпочтительно ссылаться на опубликованную версию, если она уже существует. Если вы хотите сослаться на сайт, то можно либо так же внести его в список литературы~\cite{knuthwebsite} (рекомендуется, если таких ссылок у вас много из-за особенностей темы вашего проекта), либо использовать ссылку внизу страницы\footnote{Книги доступны по ссылке: \url{http://www-cs-faculty.stanford.edu/~uno/abcde.html}, дата обр. 16.05.2013}. При работе с онлайн ресурсами не забывайте указывать дату обращения к этому ресурсу, так как в отличие от опубликованных статей, эти ресурсы могут измениться в любой момент.

% \subsection{Рисунки}

% \begin{figure}[ht]
% 	\centering
% 	\includegraphics[width=0.8\textwidth]{example.png}
% 	\caption{Пример графика. Тут должна быть подпись, поясняющая что происходит на рисунке (краткая, но достаточная для понимания основной идеи графика).}
% 	\label{fig:by_epochs}
% \end{figure}

% Все рисунки в тексте должны иметь подписи и вы на них должны ссылаться в тексте. Например, на Рисунке~\ref{fig:by_epochs} изображен пример графика. Не забывайте подписывать все оси на графиках, добавлять легенду и пояснять все обозначения, а также используйте адекватного размера шрифты и толщину линий на графиках (все должно быть видно и понятно без многократного увеличения). На рисунке из примера явно не хватает обозначения синей линии в легенде.


% \subsection{Таблицы}

% Все таблицы в тексте тоже должны иметь подписи и вы на них должны ссылаться в тексте. Например, в Таблице~\ref{table:long_epochs} показаны результаты примерного эксперимента. 


% \begin{table}[ht]
% 	\caption{Пример таблички. Тут должна быть подпись, поясняющая что происходит в таблице (краткая, но по делу).}
% 	\label{table:long_epochs}
% 	\footnotesize
% 	\centering
% 	\begin{tabular}{lrrrrrrrr}
% 		\toprule
% 		& \multicolumn{3}{c}{$\mathsf{Val}$} &
% 		\multicolumn{3}{c}{$\mathsf{Test}$} \\
% 		\cmidrule(lr){2-4} \cmidrule(l){5-7} 
% 		{} &  $\mathsf{Prec}$ &  $\mathsf{Rec}$ &  $\mathsf{F1}$ &  $\mathsf{Prec}$ &  $\mathsf{Rec}$ &  $\mathsf{F1}$  &  $\mathsf{nodes}$ & $\mathsf{subtokens}$\\
% 		\midrule
% 		запуск 1    &    0.4894 &   0.3775 &  0.4263 &     0.4824 &    0.3683 &   0.4177 & 10029 & 179\\
% 		запуск 2    &    0.4887 &   0.3739 &  0.4237 &     0.4891 &    0.3724 &   0.4228 & 10039 & 177\\
% 		запуск 3    &    0.4820 &   0.3751 &  0.4219 &     0.4838 &    0.3677 &   0.4178 & 10037&	180\\
% 		\midrule
% 		\bf{среднее} &    \bf{0.4867} &   \bf{0.3755} &  \bf{0.4239} &    \bf{ 0.4851} &    \bf{0.3695} &   \bf{0.4195} \\
% 		\bf{дисперсия}  &    0.0041 &   0.0019 &  0.0022 &     0.0036 &    0.0025 &   0.0029 \\
% 		\bottomrule
% 	\end{tabular}
% \end{table}

% \subsection{Формулы}

% Формулы стоит центрировать, а также нумеровать, если вы ссылаете на них в тексте. Также не забывайте пояснять все обозначения в формулах. Например, запишем следующую задачу оптимизации:
% \begin{equation}
%     \label{eq:si_opt}
%         \theta* = \min_{\theta} F(\theta),
% \end{equation}
% где $F$~-- квадратичная функция от параметра $\theta$. При необходимости, далее в тексте можно сослаться на формулу~(\ref{eq:si_opt}). При этом, в зависимости от конкретных формул, можно использовать разные слова: формула, уравнение, задача оптимизации и т.п.

\newpage

\section{Теоретическая часть}

\subsection{Проблема длинного контекста в больших языковых моделях}

Большие языковые модели (LLM) на основе архитектуры трансформер ~\cite{transformer} сталкиваются с фундаментальным ограничением при обработке длинных входных последовательностей. Трансформеры в настоящий момент являются одной из самых популярных архитектур для разного рода задача, в том числе и не только из области NLP. Однако квадратичная сложность операции self-attention затрудняет их эффективную работу с длинным контекстом. LLM широко применяются в разнообразных приложениях, где длина контекста значительно растёт при решении задач вроде ответов на вопросы по длинным документам, сложного логического рассуждения, агентских взаимодействий. При этом длинный контекст создаёт серьёзные трудности для эффективного инференса, включая высокие затраты памяти на KV-кэш и увеличение задержки из-за значительного объёма операций с памятью ~\cite{kv_transfer}.

Практическим следствием этих ограничений является то, что поддерживаемые контекстные окна LLM стремительно расширяются — от 4k токенов у GPT-3.5 и Llama-2 до 32k, 128k и выше. Тем не менее, даже при формально большом контекстном окне работа с контекстом длиннее того, на котором обучалась модель, может приводить к значительному снижению качества. При этом увеличивается потребление памяти и растут вычислительные затраты в связи с квадратичной сложностью self-attention.

\subsubsection{Методы сжатия контекста}

Для решения обозначенной выше проблемы существуют два подхода - инженерный и обучение латентных представлений
для сжатия. 

Первый класс — методы, не требующие дополнительного обучения и работающие с явным содержимым контекста. К ним относятся сжатие KV-кэша и prompt compression.

Ряд работ предлагает сжатие KV-кэша путём аппроксимации вычислений, однако такие методы либо безвозвратно удаляют токены, не возвращаясь к ним при последующем инференсе, либо восстанавливают предыдущие токены на уровне страниц, разделённых по текстовым позициям. Оба таких подхода ведут к потере точности выходных данных~\cite{cluster_kv}. Более семантически обоснованный подход реализован в работе ClusterKV: в ней авторы восстанавливают токены на уровне семантических кластеров и достигают пренебрежимо малых потерь точности на контекстах длиной 32k при бюджете KV-кэша всего в 1k–2k, обеспечивая при этом до 2× ускорение декодирования~\cite{cluster_kv}.

Отдельным направлением является сжатие входных данных путём удаления малоинформативных токенов. Так, в работе ~\cite{sac} предложили идею нахождения так называемых якорных токенов и сжатии всей входной информации в KV-кеши этих наиболее информативных токенов.

Второй класс методов - это методы, основанные на обучении непрерывных латентных представлений контекста в пространстве эмбеддингов. В отличие от методов, основанных на удалении малоинформативных токенов, эти  не удаляют информацию явным образом, а кодируют ее в компактное латентное представление.  

Ключевой работой в этом направлении является In-Context Autoencoder (ICAE). Его архитектура состоит из LLM-энкодера с обучаемым LoRA-адаптером, который кодирует длинный контекст в небольшое число слотов памяти, и замороженного LLM-декодера. ICAE использует датасет PwC для обучения и оценки — тот же датасет, что применяется в настоящей работе, что обеспечивает прямую сопоставимость результатов.

Другой близкой работой является данная статья ~\cite{ccot}, в которой авторы обучают вначале модель сжимать контекст в некоторые токены размышления, а затем учат модель декодер разжимать такое представление. Такой сжатый контекст может быть впоследствии использован в задачах RAG или суммаризации.

Важным ограничением таких методов является то, что при сжатии контекста длинный контекст конденсируется в небольшой набор непрерывных представлений, однако существующие методы, как правило, повторно используют сам тяжеловесный декодер в качестве обучаемого компрессора, опираясь на attention для итеративной агрегации информации. Это влечёт за собой высокие вычислительные затраты и требования к памяти за счет повтороного использования тежелого декодера.

\subsection{Cжатие цепочек рассуждений в латентном пространстве}

Отдельным направлением, непосредственно связанным с downstream-целями настоящей работы, является сжатие Chain-of-Thought (CoT) ~\cite{cot} рассуждений в непрерывное пространство. Идея такого подхода основана на том, что человек во время размышлений оперирует в глове не с конткретными словами, а с некоторыми абстракциями, концептами. Cовременные же модели рассуждений ограничены пространством человеческого языка и строят цепочки размышлений в виде последовательности дискретных токенов, представляющие фиксированные точки семантического пространства, что ограничивает их экспрессивность и приводит к неполному исследованию возможных путей рассуждения. ~\cite{soft_thinking}

В работе Soft Thinking предложен метод получения таких непрерывных токенов рассуждений и их использования: на выходе получается распределение вероятностей по словарю, затем на следующий форвард в модель уже подается линейная комбинация эмбеддингов токенов с весами из полученного распределения. Тем самым получается плавный переход от дискретных токенов в непрерывное пространство, обеспечивая тем самым более разнообразные рассуждения. ~\cite{soft_thinking}

Работа Compressed Chain-of-Thought (CCoT) исследует сжатие явных шагов рассуждения в «токены-размышления» (contemplation tokens). Метод дообучает декодерные LLM в рамках нового подхода CCoT и эмпирически оценивает их производительность и пропускную способность на бенчмарке GSM8K. Отличие CCoT подхода заключается в том, что токены-размышления подаются вместе с оставшимися дискретными токенами, а не используются просто как сигнал для декодирования. ~\cite{ccot}

\subsection{Сжатие контекста с использованием визуальной модальности}

Наиболее близким к данной работе напрвлением является использованием визуальной модальности для эффективного сжатия текстового контекста. Данный метод был предложен в работе «DeepSeek-OCR: Contexts Optical Compression»~\cite{deepseek_ocr}.

Использование визуальной модальности в качестве промежуточного пространства для сжатия текста представляет собой новое и перспективное исследовательское направление; ключевая идея состоит в том, чтобы рендерить длинный текст в одно или несколько изображений, а затем использовать мощный визуальный энкодер для сжатия высокоразмерных данных.

Работа DeepSeek-OCR представляет собой первичное исследование возможности сжатия длинных контекстов через оптическое 2D-отображение и состоит из двух компонентов: DeepEncoder и декодера DeepSeek3B-MoE-A570M. Эта работа переосмысляет vision-language модели (VLM) с LLM-центричной перспективы, фокусируясь на том, как визуальные энкодеры могут повысить эффективность LLM при обработке текстовой информации; задачи OCR, как промежуточная модальность между зрением и языком, создают естественное соответствие сжатие-декомпрессия между визуальными и текстовыми представлениями.

Вместе с тем работа «Optical Context Compression Is Just (Bad) Autoencoding» ставит под сомнение эффективность рендеринга через пиксели. DeepSeek-OCR показал, что рендеренный текст может быть восстановлен из небольшого числа визуальных токенов, но этот конвейер требует рендеринга эмбеддингов токенов в пиксели и сжатия с этого уровня — отбрасывая обученные представления в пользу изображения, которое визуальный энкодер затем должен восстановить. ~\cite{bad_ocr}. Иерархический энкодер превосходит визуальный при всех коэффициентах сжатия; для языкового моделирования визуально-основанное сжатие не превосходит даже простое усечение при сопоставимых бюджетах токенов, тогда как иерархический энкодер с этим справляется.

Настоящая работа предлагает альтернативный подход: вместо рендеринга текста в изображения, скрытые состояния текстового энкодера подаются напрямую в визуальный энкодер (ViT) через проекционный модуль. Тем самым устраняется промежуточный информационный барьер пиксельного рендеринга при сохранении архитектурной идеи использования визуального энкодера как компрессора.


\subsection{Визуальные энкодеры}

Архитектура предлагаемой модели концептуально близка к современным vision-language моделям (VLM), и в первую очередь — к BLIP-2.

BLIP-2 предлагает универсальную и эффективную стратегию предобучения, использующую готовые замороженные предобученные энкодеры изображений и замороженные текстовые декодеры. ~\cite{blip_2} Ключевым компонентом BLIP-2 является Q-Former: Q-Former — лёгкий трансформер, использующий набор обучаемых векторов-запросов для извлечения визуальных признаков из замороженного энкодера изображений; он действует как информационное узкое место между замороженным энкодером изображений и замороженным декодером. На первом этапе предобучения выполняется обучение представлениям vision-language, принуждающее Q-Former учиться визуальным представлениям, наиболее релевантным тексту; на втором этапе выполняется генеративное обучение vision-to-language путём подключения выхода Q-Former к замороженному декодеру.

Концептуально роль Q-Former в BLIP-2 аналогична роли визуального энкодера (ViT) с проекторами в предлагаемой архитектуре: оба компонента решают задачу преобразования представлений одной модальности в форму, понятную языковому декодеру. Принципиальное отличие настоящей работы состоит в инверсии направления: если BLIP-2 переводит визуальные признаки в текстовое пространство, то здесь текстовые признаки переводятся в визуальное пространство и обратно с целью сжатия.

\subsection{Сжатие цепочек рассуждений с помощью самодистилляции}

Для downstream-задач рассуждений релевантной является работа CoDi (Compressing Chain-of-Thought into Continuous Space via Self-Distillation). CoDi сжимает цепочку рассуждений в непрерывное пространство с помощью самодистилляции. ~\cite{codi} Этот подход демонстрирует, что компактные латентные представления могут сохранять достаточно информации для решения задач, требующих многошагового рассуждения, что является одной из исследовательских целей настоящей работы.




\section{Практическая часть}

\subsection{Данные}

Для обучения и валидации используется датасет sggetao/PwC, обучающая и тестовые выборки соответственно. Этот датасет использовался в обучении ICAE, поэтому это позволяет сопоставить результаты предлагаемого подхода с данным методом.

Датасет состоит из пар (промпт, длинный текстовый ответ). В рамках быстрых экспериментов в условиях ограниченных вычислительных длина контекста ограничена 1024 токенами, что позволяет оценить базовую работоспособность модели перед масштабированием.

Предобработка текста состоит в токенизации в соответствии с токенизатором, выбранного декодера, филтрации примеров по длине и формировании батчей в формате, совместимом с архитектурой модели.

\subsection{Архитектура модели}

Финальная архитектура CVLM (Compressed Vision-Language Model), используемая в практической части, представлена в Таблице~\ref{tab:architecture-final}. Общая схема обработки последовательности следующая: исходный контекст подаётся в текстовый энкодер (например, ModernBert-base), скрытые состояния которого разбиваются на непересекающиеся чанки длины $\text{cr}$ (compression rate) и сжимаются с помощью chunked attention pool в последовательность векторов длиной в $\text{cr}$ раз короче. Полученная последовательность затем проходит через MLP-проектор, визуальный энкодер (ViT) и второй MLP-проектор, приводящий размерность к скрытой размерности декодера. Сжатое представление подаётся декодеру (например, SmolLM2-1.7B-Instruct) в качестве префикса перед текстовым промптом.

Каждый компонент выбран под конкретную роль в задаче сжатия. Текстовый энкодер обеспечивает богатые контекстные представления токенов: предобучение на masked language modeling делает его подходящим для задачи понимания текста, а двунаправленный аттеншен (в отличие от каузального декодера) позволяет каждому токену учитывать весь контекст до и после, что критично, когда часть информации в дальнейшем должна быть отброшена при сжатии. Далее идет Chunked attention pool~--- механизм сжатия с обучаемыми латентными запросами: он позволяет делать обучаемое преобразование для компрессии. Визуальный энкодер, ViT встроен в пайплайн как механизм пере-агрегации компактных представлений: визуальные трансформеры обучены извлекать плотную семантику из больших массивов входных патчей, что является той же структурой, которая требуется для компрессии длинных последовательностей. Два MLP-проектора выполняют согласование размерностей между разнородными пространствами (текстовый энкодер $\to$ визуальный энкодер $\to$ декодер) и одновременно служат лёгкими обучаемыми «адаптерами» между предобученными компонентами. Декодер замыкает пайплайн: он обеспечивает обратное преобразование сжатого представления в текст и одновременно служит источником обучающего сигнала через стандартную авторегрессионную кросс-энтропию.

Ключевое преимущество такой компоновки~--- возможность \textbf{end-to-end обучения} без вспомогательных приёмов. Все компоненты~--- текстовый энкодер, attention pool, оба проектора, ViT, позиционные эмбеддинги и декодер~--- соединены дифференцируемыми операциями, поэтому градиент от единственной функции потерь (cross-entropy на токенах ответа) распространяется сквозь весь пайплайн в один прямой и обратный проход. Это устраняет необходимость в характерных для подобных задач сложностях: контрастивных вспомогательных лоссах, многостадийном предобучении, дистилляции из «учителя» с полным контекстом или обучении с подкреплением.

\begin{table}[ht]
\centering
\caption{Конфигурация архитектуры CVLM.}
\label{tab:architecture-final}
\footnotesize
\begin{tabular}{lll}
\toprule
\textbf{Компонент} & \textbf{Модель / параметры} & \textbf{Состояние при обучении}\\
\midrule
Текстовый энкодер & ModernBert-base, 22 блока, hidden=768 & Размораживается (см. \ref{sec:method-unfreeze})\\
Chunked attention pool & $K$ обучаемых латентных запросов на чанк & Обучаемый \\
Проектор 1        & MLP $768 \to 3072 \to 768$  & Обучаемый \\
Визуальный энкодер & ViT, hidden=768 & Обучаемый \\
Позиц. эмбеддинги & $V_{\max} \times 768$ для входа в ViT & Обучаемые \\
Проектор 2        & MLP $768 \to 3072 \to 2048$ & Обучаемый \\
Языковой декодер  & SmolLM2-1.7B-Instruct & Обучаемый \\
\bottomrule
\end{tabular}
\end{table}

В качестве декодера выбрана SmolLM2-1.7B-Instruct: она использует Grouped-Query Attention (GQA)~\cite{gqa}, обладает достаточной ёмкостью для содержательной оценки качества восстановления и обеспечивает приемлемую скорость итераций обучения на одной GPU класса A100/H100. Текстовый энкодер ModernBert-base изначально замораживался полностью; последующие эксперименты показали, что его разморозка является ключевым фактором улучшения качества (раздел~\ref{sec:exp-unfreeze}). В CVLM декодер обучается вместе с остальными компонентами; замороженный декодер используется только в оценочных бейзлайнах \texttt{baseline\_llm} и \texttt{baseline\_llm\_full}.

\subsubsection{Chunked attention pool как информационное «бутылочное горлышко»}\label{sec:method-pool}

Сжатие реализовано как пулинг внутри непересекающихся чанков длины $\text{cr}$. Для последовательности скрытых состояний $H \in \mathbb{R}^{L \times d}$ после разбиения на чанки $C_j \in \mathbb{R}^{\text{cr} \times d}$ ($j = 1, \dots, \lceil L/\text{cr} \rceil$) каждый чанк обрабатывается attention pool:
\begin{equation}
    \label{eq:attn-pool}
    o_j = \frac{1}{K} \sum_{k=1}^{K} \mathrm{softmax}\!\left( \frac{C_j q_k^\top}{\sqrt{d}} \right)^\top C_j,
\end{equation}
где $q_k \in \mathbb{R}^{d}$~--- обучаемые латентные запросы ($k = 1, \dots, K$), а маски паддинга учитываются установкой $-\infty$ в соответствующих позициях до softmax. При $K = 1$ выражение~(\ref{eq:attn-pool}) соответствует обычному attention pool с одним запросом, при $K > 1$ становится более общим видом пулинга. Существенно, что при нулевой инициализации запросов $q_k = 0$ softmax становится равномерным и формула вырождается в обычный пулинг (mean-pool) для любого $K$. Это обеспечивает строгую обратную совместимость с базовым подходом и позволяет модели выучивать более сложные зависимости.

\subsubsection{Куррикулум по коэффициенту сжатия}\label{sec:method-curriculum}

Прямое обучение модели при высоком $\text{cr}$ ведёт к плохой сходимости, поскольку текстовый энкодер, визуальный энкодер и текстовый декодер инициализрованы весами, которые явным образом не обучались на задачу сжатия, при этом визуальный энкодер не получал преобразованные текстовые эмбеддинги ранее, поэтому обучение может проходить нестабильно. Для решения этой проблемы реализован куррикулум обучения по $\text{cr}$. Расписание задаётся последовательностью пар $(\text{cr}_i, t_i)$, где $t_i$~--- глобальный шаг оптимизатора, после которого происходит переход к стадии $i+1$. При смене стадии состояние оптимизатора (моменты Адама) переносится без сброса, что обеспечивает плавное продолжение обучения; на каждом переходе принудительно сохраняется чекпоинт для последующей покомпонентной оценки.

В экспериментах используются два расписания: $1 \to 2 \to 4 \to 8$ (4-стадийный, 8 эпох) и $1 \to 2 \to 4 \to 8 \to 16$ (5-стадийный, 10 эпох). Численное сравнение представлено в разделе~\ref{sec:exp-cr16}.

\subsubsection{Разморозка текстового энкодера}\label{sec:method-unfreeze}

Полностью замороженный текстовый энкодер представляет собой существенное узкое место: его представления оптимизированы под задачу masked language modeling, но не под задачу сжатия в визуальное латентное пространство. Поэтому в экспериментах был добавлен механизм разморозки верхних $K$ блоков текстового энкодера; при $K = 22$ (полное число блоков в ModernBert-base) энкодер становится полностью обучаемым.

\subsection{Этапы обучения}

Этап 1 заключается в обучении на реконструкцию - модель обучается восстанавливать исходный текст по его сжатому латентному представлению. На данном этапе функцией потерь является кросс-энтропия на токенах восстановленного текста.

На втором этапе планируется тонкая донастройка под downstream-задачи, то есть после достижения приемлемого качества реконструкции будет исследовано применение сжатых представлений в задачах, требующих рассуждений (reasoning).

\subsection{Метрики качества}\label{sec:metrics}

В работе используются метрики двух типов: \emph{качество восстановления} (насколько хорошо модель предсказывает целевой ответ при условии сжатого контекста) и \emph{эффективность сжатия} (насколько компактным оказывается это представление по отношению к исходному документу). Обозначения, единые для всего раздела: $S_i$~--- число токенов исходного документа в токенизаторе декодера для примера $i$, $V_i$~--- число vision-токенов после chunked attention pool, $P_i$~--- число токенов промпта (вопрос + системные токены), $A_i$~--- число токенов референсного ответа, $N_{\text{ans}} = \sum_i |A_i|$~--- суммарное число токенов ответов на тестсете. Через $\mathrm{NLL}(A_i)$ обозначается суммарный кросс-энтропийный логарифм правдоподобия (в натах) на токенах ответа примера $i$ при teacher forcing:
\begin{equation}
\label{eq:nll}
    \mathrm{NLL}(A_i) = -\sum_{t=1}^{|A_i|} \ln p_\theta\!\left(a_{i,t} \mid a_{i,<t},\, \text{ctx}_i\right),
\end{equation}
где $\text{ctx}_i$~--- сжатый контекст (vision-токены + промпт) для примера $i$.

\paragraph{Perplexity (PPL).} Стандартная мера «удивлённости» модели на токенах ответа:
\begin{equation}
\label{eq:ppl}
    \mathrm{PPL} = \exp\!\left( \frac{1}{N_{\text{ans}}} \sum_{i} \mathrm{NLL}(A_i) \right).
\end{equation}
Ниже~--- лучше; интерпретируется как эффективный размер словаря, между токенами которого модель «колеблется» в среднем.

\paragraph{Token accuracy (TA).} Доля токенов ответа, для которых $\arg\max$ предсказательного распределения совпадает с целевым токеном:
\begin{equation}
\label{eq:tokacc}
    \mathrm{TA} = \frac{1}{N_{\text{ans}}} \sum_{i} \sum_{t=1}^{|A_i|} \mathbb{1}\!\left[\arg\max_{v} p_\theta\!\left(v \mid a_{i,<t},\, \text{ctx}_i\right) = a_{i,t}\right].
\end{equation}
Выше~--- лучше. В отличие от PPL, не зависит от калибровки распределения, что делает TA устойчивой при сравнении моделей с разной плотностью предсказаний.


\paragraph{Compression ratio (CR).} Отношение длины исходного документа в токенах декодера к длине его сжатого представления:
\begin{equation}
\label{eq:cr}
    \mathrm{CR}_i = \frac{S_i}{V_i}, \qquad
    \overline{\mathrm{CR}} = \frac{1}{N} \sum_{i=1}^{N} \frac{S_i}{V_i}.
\end{equation}
Выше~--- сильнее сжатие. Эмпирическое распределение $\{\mathrm{CR}_i\}$ может отличаться от номинального $\text{cr}$-гиперпараметра, поскольку $S_i$ считается в токенизаторе декодера, а $V_i$ выводится из токенизации энкодера; на практике для $\text{cr}=8$ среднее $\overline{\mathrm{CR}} \approx 7$. В работе сообщаются среднее, медиана и перцентили $\mathrm{p}_{10}/\mathrm{p}_{90}$.

\paragraph{Bits per source token (BPST).} Совместная метрика качества и сжатия: кросс-энтропия на ответе, выраженная в битах и отнесённая к суммарному числу токенов \emph{исходного} документа:
\begin{equation}
\label{eq:bpst}
    \mathrm{BPST} = \frac{\sum_{i} \mathrm{NLL}(A_i) \,/\, \ln 2}{\sum_{i} S_i}.
\end{equation}
Ниже~--- лучше. В отличие от PPL, эта величина \emph{сравнима между разными $\text{cr}$}: PPL зависит только от длины ответа, тогда как BPST явно учитывает, сколько токенов исходного контекста стояло за каждым предсказанным битом ответа.

\paragraph{Effective context reduction (ECR).} Эффективный коэффициент сжатия по отношению к длине контекста, реально потребляемого декодером:
\begin{equation}
\label{eq:ecr}
    \mathrm{ECR} = \frac{\overline{S}}{\overline{V} + \overline{P}},
\end{equation}
где $\overline{S}$, $\overline{V}$, $\overline{P}$~--- средние по тестсету $S_i$, $V_i$, $P_i$ соответственно. $\mathrm{ECR} > 1$ означает, что CVLM подаёт декодеру в $\mathrm{ECR}$ раз меньше токенов, чем потребовалось бы для полнотекстового бейзлайна; $\mathrm{ECR} = 1$~--- безубыточный режим; $\mathrm{ECR} < 1$~--- сжатое представление длиннее исходного (вырожденный случай). Для бейзлайнов \texttt{baseline\_llm\_full} и SFT знаменатель заменяется на среднюю длину полнотекстового промпта $\overline{P}_{\text{full}}$, для \texttt{baseline\_llm} метрика не определена ($\overline{P} = 0$ по документу).

\paragraph{Сводка.} Метрики разделяются на три группы по интерпретации (Таблица~\ref{tab:metrics-summary}).

\begin{table}[ht]
\centering
\caption{Метрики, используемые в работе.}
\label{tab:metrics-summary}
\footnotesize
\begin{tabular}{lll}
\toprule
\textbf{Группа} & \textbf{Метрика} & \textbf{Формула} \\
\midrule
\multirow{2}{*}{Качество восстановления} & Perplexity (PPL) $\downarrow$  & \eqref{eq:ppl} \\
                                                  & Token accuracy (TA) $\uparrow$ & \eqref{eq:tokacc} \\
\midrule
\multirow{3}{*}{Эффективность сжатия}             & Compression ratio (CR) $\uparrow$   & \eqref{eq:cr} \\
                                                  & Bits per source token (BPST) $\downarrow$ & \eqref{eq:bpst} \\
                                                  & Effective context reduction (ECR) $\uparrow$ & \eqref{eq:ecr} \\
\bottomrule
\end{tabular}
\end{table}

\subsection{Ожидаемые результаты и анализ}

Полный план экспериментов охватывает:

\begin{itemize}
    \item Верификацию базовой работоспособности модели~--- уверенное восстановление текста с приемлемой точностью (\textbf{выполнено}, см.~раздел~\ref{sec:experiments});
    \item Сравнение качества реконструкции с SFT-бейзлайном и ICAE (см.~раздел~\ref{sec:future});
    \item Исследование зависимости качества от типа визуального энкодера (см.~раздел~\ref{sec:future});
    \item Исследование зависимости качества от типа текстового энкодера (см.~раздел~\ref{sec:future});
    \item Исследование зависимости качества реконструкции от коэффициента сжатия (\textbf{первичные результаты получены}, см.~разделы~\ref{sec:exp-cr16},~\ref{sec:experiments-summary});
    \item Оценку применимости сжатых представлений в downstream-задачах на рассуждения (см.~раздел~\ref{sec:future});
    \item Анализ вычислительной эффективности предложенного подхода.
\end{itemize}

По итогам ожидается определить оптимальную конфигурацию архитектуры, позволяющую сохранять приемлемое качество при сжатии контекста, и далее использовать её в downstream-задачах, в первую очередь требующих длинных цепочек рассуждений.

\section{Реализация и проведённые эксперименты}\label{sec:experiments}

В данном разделе представлены результаты практической части работы. Реализованная модель CVLM прошла серию итеративных улучшений, каждое из которых описано в виде отдельного эксперимента с количественным сравнением. Все эксперименты выполнены на одной GPU класса A100/H100, обучение в bf16 с активированным gradient checkpointing для декодера и (по необходимости) текстового энкодера.

\subsection{Общая методология оценки}\label{sec:eval-methodology}

Обучение и оценка проводятся на датасете \texttt{sggetao/PwC} (training и test split соответственно). Целевые метрики:

\begin{itemize}
    \item \textbf{Perplexity (PPL)} $\downarrow$~--- считается на токенах ground-truth ответа в teacher forcing-режиме, $\mathrm{PPL} = \exp(\mathrm{CE\text{-}loss})$;
    \item \textbf{ROUGE-L} $\uparrow$~--- F-мера на основе наибольшей общей подпоследовательности, считается на сгенерированных моделью ответах;
    \item \textbf{BLEU-4} $\uparrow$~--- стандартная n-граммовая мера качества генерации;
    \item \textbf{Token accuracy} $\uparrow$~--- доля токенов ground-truth ответа, для которых argmax распределения декодера в teacher forcing-режиме совпадает с истинным токеном. В отличие от PPL, метрика отражает не калибровку распределения, а его top-1 точность; в отличие от ROUGE-L, считается на каждом токене независимо, без необходимости генерации. Сравнима между запусками за счёт идентичного тестового сплита.
\end{itemize}

Для каждого запуска сохраняются промежуточные чекпоинты на границах стадий куррикулума. Каждый чекпоинт оценивается \textit{в нативном режиме}~--- при том $\text{cr}$, при котором веса обучались, что устраняет систематический сдвиг при сравнении мульти-стадийных запусков.

В качестве референсных границ используются два бейзлайна на той же декодерной модели:
\begin{itemize}
    \item \texttt{baseline\_llm} (без документа, только промпт)~--- нижняя граница, ROUGE-L~$\approx 0.097$;
    \item \texttt{baseline\_llm\_full} (полный документ без сжатия, оракул)~--- верхняя граница, ROUGE-L~$\approx 0.165$--$0.168$.
\end{itemize}
Эти значения стабильны между запусками (различие $< 0.005$), что подтверждает корректность пайплайна оценки.

\subsection{Эксперимент 1: от mean-pool к attention pool}\label{sec:exp-pool}

В исходной реализации в качестве операции сжатия использовался средний пулинг (mean-pool) внутри каждого чанка длины $\text{cr}$. Этот подход не имеет обучаемых параметров и не способен выделять более информативные токены внутри чанка. На его место предложен chunked attention pool с одним обучаемым латентным запросом ($K = 1$), инициализируемым нулями. Эта инициализация гарантирует бит-эквивалентность mean-pool на нулевом шаге обучения, обеспечивая контролируемый запуск.

Эксперимент проведён при идентичных всех остальных параметрах. Результаты на финальной стадии куррикулума $\text{cr}=8$ (шаг 30201) представлены в Таблице~\ref{tab:exp-pool}.

\begin{table}[ht]
\centering
\caption{Сравнение методов пулинга при идентичной конфигурации (замороженный энкодер, $\text{cr}=8$, шаг~30201).}
\label{tab:exp-pool}
\begin{tabular}{lccc}
\toprule
\textbf{Метод пулинга} & \textbf{PPL} $\downarrow$ & \textbf{ROUGE-L} $\uparrow$ & \textbf{Token acc.} $\uparrow$ \\
\midrule
Mean-pool (baseline)      & 6.13 & 0.111 & 0.591 \\
\textbf{Attention pool ($K=1$)}      & \textbf{5.50} & \textbf{0.116} & \textbf{0.605} \\
\bottomrule
\end{tabular}
\end{table}

Замена среднего пулинга на attention pool с одним обучаемым латентом дала относительное снижение PPL на $10.3\%$, рост ROUGE-L на $4.5\%$ и рост token accuracy на $+1.4$ п.п. Это подтверждает гипотезу о пользе обучаемого механизма выбора информативных токенов внутри чанка и одновременно показывает, что даже минимальное число обучаемых параметров пулинга существенно влияет на качество.

\subsection{Эксперимент 2: разморозка текстового энкодера}\label{sec:exp-unfreeze}

В исходной конфигурации ModernBert-base был полностью заморожен ради экономии памяти и сохранения предобученных представлений. Анализ промежуточных представлений показал, что замороженный энкодер не способен адаптироваться к задаче: исходные эмбеддинги токенов оптимизированы под MLM-задачу, но не под последующее сжатие в визуальное латентное пространство.

Реализован параметризуемый механизм разморозки верхних $K$ блоков энкодера с \texttt{final\_norm}; в данном эксперименте $K = 22$ (полное число блоков). Все остальные компоненты идентичны Эксперименту~\ref{sec:exp-pool}. Результаты на финальной стадии $\text{cr}=8$ представлены в Таблице~\ref{tab:exp-unfreeze}.

\begin{table}[ht]
\centering
\caption{Влияние разморозки текстового энкодера ($\text{cr}=8$, шаг~30201). Полная разморозка $K=22$ существенно сокращает разрыв до верхней границы оракула.}
\label{tab:exp-unfreeze}
\begin{tabular}{lccc}
\toprule
\textbf{Конфигурация} & \textbf{PPL} $\downarrow$ & \textbf{ROUGE-L} $\uparrow$ & \textbf{Token acc.} $\uparrow$ \\
\midrule
Замороженный энкодер, mean-pool           & 6.13 & 0.111 & 0.591 \\
Замороженный энкодер, attention pool      & 5.50 & 0.116 & 0.605 \\
\textbf{Полностью размороженный (K=22), attention pool} & \textbf{3.77} & \textbf{0.133} & \textbf{0.666} \\
\midrule
\textit{baseline\_llm} (без документа, нижняя граница) & --- & 0.097 & --- \\
\textit{baseline\_llm\_full} (полный документ, оракул) & --- & 0.165 & --- \\
\bottomrule
\end{tabular}
\end{table}

Это основной количественный результат первой части работы:
\begin{itemize}
    \item PPL снижается на $38.5\%$ относительно базовой конфигурации mean-pool (с 6.13 до 3.77);
    \item ROUGE-L вырастает на $19.8\%$ (с 0.111 до 0.133);
    \item Token accuracy вырастает на $+7.5$ п.п. (с $0.591$ до $0.666$), т.е. модель верно предсказывает следующий токен в 2/3 случаев при сжатии контекста в 8 раз;
    \item Разрыв до верхнего оракула сокращается с $33\%$ до $20\%$ относительных пунктов ROUGE-L, при этом сжатие контекста сохраняется в 8 раз.
\end{itemize}

Это подтверждает, что текстовый энкодер действительно является ключевым узким местом, и его адаптация к задаче сжатия даёт качественный скачок, недостижимый ни увеличением числа эпох, ни усложнением пулинга при замороженном энкодере.

\subsection{Эксперимент 3: расширение куррикулума до $\text{cr}=16$}\label{sec:exp-cr16}

Естественным продолжением является вопрос о возможности сжатия в 16 раз. Был проведён эксперимент с расширенным куррикулумом $1 \to 2 \to 4 \to 8 \to 16$ (расписание $1{:}6000, 2{:}12000, 4{:}18000, 8{:}24000, 16{:}0$) на конфигурации из Эксперимента~\ref{sec:exp-unfreeze} (полностью размороженный энкодер, attention pool $K=1$, 10 эпох). Дополнительно представлены результаты двух связанных запусков: укороченного 4-стадийного куррикулума (8 эпох, обозначен d1) и 4-стадийного на 10 эпохах с длительной фазой $\text{cr}=8$, оцененного также при $\text{cr}=16$ в режиме zero-shot экстраполяции (обозначен d2).

Сравнение по стадиям представлено в Таблице~\ref{tab:exp-cr16}.

\begin{table}[ht]
\centering
\caption{Динамика PPL и ROUGE-L по стадиям куррикулума. d1 - 4-стадийный куррикулум, 8 эпох, остановка на $\text{cr}=8$. d3 - 5-стадийный куррикулум, 10 эпох, последняя стадия $\text{cr}=16$. d2 - 4-стадийный куррикулум на 10 эпох (длинная фаза $\text{cr}=8$), оценка финальных весов при $\text{cr}=16$ в zero-shot режиме. Звёздочка ($^{*}$) обозначает оценку с экстраполяцией к $\text{cr}=16$.}
\label{tab:exp-cr16}
\footnotesize
\begin{tabular}{lcccccc}
\toprule
& \multicolumn{2}{c}{d1 (8 эпох, cr$\to$8)} & \multicolumn{2}{c}{d3 (10 эпох, cr$\to$16)} & \multicolumn{2}{c}{d2 (10 эпох, eval@$\text{cr}=16$)} \\
\cmidrule(lr){2-3}\cmidrule(lr){4-5}\cmidrule(lr){6-7}
Шаг & PPL & R-L & PPL & R-L & PPL & R-L \\
\midrule
12000 ($\text{cr}=2$)           & 2.80           & 0.122          & 2.66           & 0.121          & 2.60           & 0.124 \\
18000 ($\text{cr}=4$)           & 3.17           & \textbf{0.136} & 3.59           & 0.131          & 3.06           & \textbf{0.141} \\
22651 ($\text{cr}=8$)           & 3.88           & 0.129          & 4.70           & 0.120          & 3.80           & 0.131 \\
30201 ($\text{cr}=8$ для d1)    & \textbf{3.77}  & \textbf{0.133} & 5.64$^{*}$     & 0.115$^{*}$    & 5.23$^{*}$     & 0.120$^{*}$ \\
37751 (конец)                   & ---            & ---            & 5.63           & 0.116          & 5.22$^{*}$     & 0.120$^{*}$ \\
\bottomrule
\end{tabular}
\end{table}

Для большей наглядности динамика top-1 точности декодера по стадиям куррикулума вынесена в отдельную Таблицу~\ref{tab:exp-cr16-tokacc}. Здесь явно видно, что d2 (длительная фаза $\text{cr}=8$) удерживает более высокую token accuracy на всех стадиях, кроме явного обучения на $\text{cr}=16$, где он оценивается в режиме zero-shot.

\begin{table}[ht]
\centering
\caption{Token accuracy по стадиям куррикулума для тех же трёх запусков (d1, d3, d2). Жирным выделено максимальное значение в строке; звёздочкой~--- оценка при $\text{cr}=16$ без явного обучения на этой стадии.}
\label{tab:exp-cr16-tokacc}
\footnotesize
\begin{tabular}{lccc}
\toprule
Шаг (стадия $\text{cr}$) & \textbf{d1} (8 эпох) & \textbf{d3} (5-стадий, cr$\to$16) & \textbf{d2} (4-стадий, eval@$\text{cr}=16$) \\
\midrule
6000 ($\text{cr}=1$)    & 0.574          & \textbf{0.768} & 0.746 \\
12000 ($\text{cr}=2$)   & 0.726          & 0.737          & \textbf{0.742} \\
18000 ($\text{cr}=4$)   & 0.700          & 0.677          & \textbf{0.707} \\
22651 ($\text{cr}=8$)   & 0.659          & 0.629          & \textbf{0.663} \\
30201 ($\text{cr}=8 / \text{cr}=16$ для d2,d3) & \textbf{0.666} & 0.599 & 0.610$^{*}$ \\
37751 (конец)           & ---            & 0.601          & 0.611$^{*}$ \\
\bottomrule
\end{tabular}
\end{table}

Из таблиц следуют два важных наблюдения, влияющих на дальнейший план экспериментов:

\textbf{Наблюдение 1.} $\text{cr}=4$ является эмпирическим оптимумом по качеству генерации во всех конфигурациях: ROUGE-L достигает $0.131$--$0.141$ против $0.122$--$0.133$ на $\text{cr}=8$ и $0.115$--$0.120$ на $\text{cr}=16$. Это даёт операционно полезный режим со сжатием в 4 раза и просадкой $\sim 0.025$ ROUGE-L относительно оракула.

\textbf{Наблюдение 2.} Явное обучение на $\text{cr}=16$ хуже, чем нулевая экстраполяция с $\text{cr}=8$, при равном бюджете эпох. В d3 модель обучалась 13750 шагов при $\text{cr}=16$ и пришла к PPL $5.63$ / R-L $0.116$. В d2 модель обучалась $\sim 20000$ шагов на $\text{cr}=8$ и оценивалась на $\text{cr}=16$ без какого-либо обучения на этом сжатии~--- PPL $5.22$ / R-L $0.120$, лучше по обоим показателям. Объяснение состоит в том, что расширение куррикулума до $\text{cr}=16$ "размывает" бюджет обучения и оставляет $\text{cr}=8$ существенно недообученным (6000 шагов в d3 против 12000 в d1). Это, в свою очередь, ухудшает базу, с которой модель переходит к $\text{cr}=16$. Из этого следует прямое практическое следствие: расписание куррикулума требует пересмотра в пользу более длительной фазы $\text{cr}=8$ (см. направление~\ref{todo:reschedule}).

\begin{figure}[ht]
    \centering
    \includegraphics[width=0.9\textwidth]{training_curves_K4.pdf}
    \caption{Динамика PPL, ROUGE-L и token accuracy по стадиям куррикулума для лучшей конфигурации $K=4$ (d4). Фон показывает текущий коэффициент сжатия.}
    \label{fig:training-curves-k4}
\end{figure}

\subsection{Эксперимент 4: multi-query attention pool}\label{sec:exp-multilatent}

Логичным расширением attention pool с одним латентом является использование $K > 1$ обучаемых латентных запросов на чанк по аналогии с подходом Perceiver. В реализации каждый чанк независимо опрашивается всеми $K$ латентами по формуле~(\ref{eq:attn-pool}), и затем $K$ полученных представлений усредняются в один вектор. Геометрия выходной последовательности не меняется (по-прежнему один токен на чанк), и сохраняется строгая обратная совместимость: при нулевой инициализации запросов любая $K$ снова бит-эквивалентна mean-pool. Реализация покрыта четырьмя unit-тестами, проверяющими: (i) совпадение с mean-pool при нулевой инициализации для $K \in \{1, 4\}$; (ii) отклонение от mean-pool при ненулевой инициализации; (iii) сохранение размерности выхода для $K \in \{1, 2, 4, 8\}$.

Эксперимент проведён при $K = 4$ и идентичных всех остальных параметрах относительно d3 (полная разморозка энкодера, расписание $1 \to 2 \to 4 \to 8 \to 16$, 10 эпох). Этот запуск обозначен как d4. Тренировочная кросс-энтропия снизилась с $1.124$ ($K=1$, d3) до $1.104$ ($K=4$). Однако полная оценка на тестовой выборке показывает существенно больший эффект, чем предсказывал тренировочный лосс: $K=4$ доминирует над $K=1$ на всех стадиях сжатия $\text{cr} \in \{2, 4, 8, 16\}$ по всем трём метрикам. Сравнение по стадиям приведено в Таблице~\ref{tab:exp-multilatent}.

\begin{table}[ht]
\centering
\caption{Сравнение $K=4$ (d4) и $K=1$ (d3) при идентичной конфигурации (полностью размороженный энкодер, 5-стадийный куррикулум $1 \to 2 \to 4 \to 8 \to 16$, 10 эпох). Нативная оценка~--- каждый чекпоинт оценивается при том $\text{cr}$, при котором он обучался. Жирным выделен лучший результат в строке.}
\label{tab:exp-multilatent}
\footnotesize
\begin{tabular}{lccccccc}
\toprule
& & \multicolumn{3}{c}{$K=1$ (d3)} & \multicolumn{3}{c}{$K=4$ (d4)} \\
\cmidrule(lr){3-5}\cmidrule(lr){6-8}
Шаг & cr & PPL$\downarrow$ & R-L$\uparrow$ & TokAcc$\uparrow$ & PPL$\downarrow$ & R-L$\uparrow$ & TokAcc$\uparrow$ \\
\midrule
6000  & 1   & \textbf{2.31} & 0.104          & \textbf{0.768} & 2.41          & \textbf{0.104} & 0.759 \\
12000 & 2   & 2.66          & 0.121          & 0.737          & \textbf{2.56} & \textbf{0.125} & \textbf{0.745} \\
18000 & 4   & 3.59          & 0.131          & 0.677          & \textbf{3.00} & \textbf{0.142} & \textbf{0.711} \\
24000 & 8   & 4.66          & 0.123          & 0.631          & \textbf{3.68} & \textbf{0.135} & \textbf{0.670} \\
37751 & 16  & 5.63          & 0.116          & 0.601          & \textbf{4.45} & \textbf{0.128} & \textbf{0.635} \\
\bottomrule
\end{tabular}
\end{table}

Это второй основной количественный результат работы. На $\text{cr}=8$ PPL снижается с $4.66$ до $3.68$ ($-21\%$), ROUGE-L растёт с $0.123$ до $0.135$ ($+9.8\%$), token accuracy~--- с $0.631$ до $0.670$ ($+3.9$ п.п.). Сопоставимое относительное улучшение наблюдается и на $\text{cr}=16$: PPL $-21\%$, R-L $+10.2\%$. На малых $\text{cr} \in \{1, 2\}$ преимущество $K=4$ остаётся, но численно меньше, что согласуется с тем, что при низком сжатии каждый чанк короче и роль механизма пулинга ослабевает.

Более того, лучшая по совокупности метрика при $\text{cr}=8$ достигается именно d4 на шаге 24000: PPL $3.68$ / R-L $0.135$ / TokAcc $0.670$. Это превосходит лучший показатель d1 ($K=1$, 4-стадийный куррикулум), который был зафиксирован как основной результат Эксперимента~\ref{sec:exp-unfreeze}: PPL $3.77$ / R-L $0.133$ / TokAcc $0.666$. Таким образом, $K=4$ при 5-стадийном куррикулуме одновременно (i)~побивает прежний лучший результат на $\text{cr}=8$ и (ii)~открывает практически применимый режим $\text{cr}=16$ с ROUGE-L $0.128$ (то есть выше нижней границы бейзлайна $0.097$).

Интерпретация: разрыв между приростом тренировочной CE ($-1.8\%$) и приростом на тестовых метриках ($-21\%$ PPL) говорит о том, что мульти-латентный пулинг улучшает не «среднюю» точность распределения декодера на легких токенах (которые доминируют в усреднении CE), а извлечение трудных для восстановления токенов~--- что отражается именно в PPL и ROUGE-L. Содержательно это согласуется с гипотезой о том, что $K$ независимых латентов извлекают несколько ортогональных семантических осей внутри одного чанка, что критично для длинных документов с гетерогенным содержимым (имена, цифры, технические термины).

\subsection{Сводка результатов}\label{sec:experiments-summary}

Совокупный прогресс первой части работы относительно исходного варианта представлен в Таблице~\ref{tab:summary}.

\begin{table}[ht]
\centering
\caption{Сводка прогресса первой части работы. Основные строки~--- при $\text{cr}=8$, нативная оценка. Дополнительно приведены лучший режим по качеству ($\text{cr}=4$) и практический режим высокого сжатия ($\text{cr}=16$).}
\label{tab:summary}
\begin{tabular}{lccccc}
\toprule
\textbf{Конфигурация} & \textbf{PPL} $\downarrow$ & \textbf{ROUGE-L} $\uparrow$ & \textbf{Token acc.} $\uparrow$ & $\Delta$ R-L & $\Delta$ TA \\
\midrule
\multicolumn{6}{l}{\emph{Основной режим $\text{cr}=8$}} \\
Mean-pool, замороженный энкодер              & 6.13          & 0.111          & 0.591          & 0.000           & 0.000 \\
\;\;+ attention pool ($K=1$)                 & 5.50          & 0.116          & 0.605          & +0.005          & +0.014 \\
\;\;+ полная разморозка ($K=22$, 4-стадии)   & 3.77          & 0.133          & 0.666          & +0.022          & +0.075 \\
\;\;+ multi-query pool ($K{=}4$, 5-стадии)   & \textbf{3.68} & \textbf{0.135} & \textbf{0.670} & \textbf{+0.024} & \textbf{+0.079} \\
\midrule
\multicolumn{6}{l}{\emph{Дополнительно}} \\
$\text{cr}=4$ ($K{=}4$, тот же запуск)       & 3.00          & 0.142          & 0.711          & +0.031          & +0.120 \\
$\text{cr}=16$ ($K{=}4$, обучение)           & 4.45          & 0.128          & 0.635          & +0.017          & +0.044 \\
\bottomrule
\end{tabular}
\end{table}

\begin{figure}[ht]
    \centering
    \includegraphics[width=\textwidth]{progression_summary.pdf}
    \caption{Совокупный вклад ключевых изменений архитектуры при $\text{cr}=8$: переход от mean-pool к attention pool, полная разморозка текстового энкодера и multi-query attention pool с $K=4$.}
    \label{fig:progression-summary}
\end{figure}

\begin{figure}[ht]
    \centering
    \includegraphics[width=0.78\textwidth]{pareto_cr_rougeL.pdf}
    \caption{Парето-кривая качества восстановления: ROUGE-L в зависимости от коэффициента сжатия для $K=1$ и $K=4$ с нижней и верхней границами качества.}
    \label{fig:pareto-cr-rougel}
\end{figure}

\begin{table}[ht]
\centering
\caption{Качество и эффективность сжатия для лучшей конфигурации $K=4$ (d4). Чекпоинт выбирается по максимальному ROUGE-L на соответствующей стадии $\text{cr}$.}
\label{tab:compression-efficiency}
\footnotesize
\begin{tabular}{cccccccc}
\toprule
$\text{cr}$ & Шаг & PPL$\downarrow$ & R-L$\uparrow$ & TokAcc$\uparrow$ & $\overline{\mathrm{CR}}$ & BPST$\downarrow$ & ECR$\uparrow$ \\
\midrule
1 & 6000 & 2.41 & 0.104 & 0.759 & 1.05 & 0.113 & 1.03 \\
2 & 11326 & 2.55 & 0.125 & 0.745 & 2.10 & 0.120 & 2.00 \\
4 & 18000 & 3.00 & 0.142 & 0.711 & 4.19 & 0.141 & 3.81 \\
8 & 24000 & 3.68 & 0.135 & 0.670 & 8.36 & 0.168 & 6.96 \\
16 & 33976 & 4.46 & 0.128 & 0.634 & 16.59 & 0.192 & 11.85 \\
\bottomrule
\end{tabular}
\end{table}


Таблица~\ref{tab:compression-efficiency} дополняет метрики качества явными метриками эффективности. Для основного режима $\text{cr}=8$ эмпирическое среднее отношение сжатия составляет $8.36$, а эффективное сокращение контекста, реально подаваемого в декодер с учётом промпта, равно $6.96$. Поэтому заявленное сжатие не является только номинальным гиперпараметром: оно отражается в фактической длине входа декодера.

Ключевые количественные выводы:
\begin{enumerate}
    \item Совокупный прогресс от трёх нововведений (attention pool, полная разморозка, multi-query $K=4$): \textbf{PPL $-40.0\%$, ROUGE-L $+21.6\%$, token accuracy $+13.4\%$ (с $0.591$ до $0.670$)} при том же сжатии в 8 раз.
    \item Из них два независимых вклада~--- разморозка энкодера ($K=22$) и мульти-латентный пулинг ($K=4$)~--- складываются почти аддитивно по ROUGE-L: разморозка даёт $+0.017$ (с $0.116$ до $0.133$), мульти-латентный пулинг~--- ещё $+0.002$ при тех же остальных компонентах и $-21\%$ PPL.
    \item Разрыв до верхнего оракула (полный документ без сжатия, R-L $\approx 0.165$) сокращён с $33\%$ до $18\%$ относительных пунктов ROUGE-L.
    \item Оптимальный рабочий режим по соотношению качество/сжатие: $\text{cr}=4$ при $K=4$~--- ROUGE-L $0.142$ при сжатии в 4 раза, что близко к верхнему оракулу.
    \item Режим $\text{cr}=16$ становится практически применимым благодаря $K=4$: ROUGE-L $0.128$ (выше нижней границы бейзлайна $0.097$ на $32\%$) при сжатии в 16 раз. На $K=1$ при том же расписании режим $\text{cr}=16$ деградирует до $0.116$.
\end{enumerate}

\subsection{Сравнение с бейзлайнами}\label{sec:baselines}

Для интерпретации абсолютных значений метрик необходимо позиционировать каждую конфигурацию между двумя референсными точками, оценёнными на той же декодерной модели и тестовом сплите:

\begin{itemize}
    \item \texttt{baseline\_llm}~--- декодер видит только вопрос без какого-либо документа. Это нижняя граница: качество, достижимое исключительно за счёт параметрической памяти декодера.
    \item \texttt{baseline\_llm\_full}~--- декодер видит полный документ без сжатия (oracle). Это верхняя граница, к которой стремятся методы сжатия.
\end{itemize}

Все CVLM-конфигурации должны попадать в коридор между этими двумя точками; чем ближе к верхней границе при большем сжатии~--- тем сильнее метод.

\paragraph{Сравнение всех конфигураций при cr=8.} В Таблице~\ref{tab:vs-baselines-cr8} сведены все обученные конфигурации работы при фиксированном сжатии в 8 раз вместе с обеими реперами. Колонки «закрыт разрыв» показывают долю расстояния между нижней и верхней границей, пройденную данной конфигурацией, по token accuracy и ROUGE-L соответственно.

\begin{table}[ht]
\centering
\caption{Полное сравнение конфигураций при $\text{cr}=8$ с обоими бейзлайнами. «Закрыт разрыв (TA)» вычисляется как $(\mathrm{TA} - \mathrm{TA}_{\text{lower}}) / (\mathrm{TA}_{\text{upper}} - \mathrm{TA}_{\text{lower}})$, аналогично для R-L. $100\%$~--- достижение оракула, $0\%$~--- неотличимо от no-document бейзлайна.}
\label{tab:vs-baselines-cr8}
\footnotesize
\begin{tabular}{lcccccc}
\toprule
\textbf{Конфигурация} & \textbf{PPL} $\downarrow$ & \textbf{TokAcc} $\uparrow$ & \textbf{R-L} $\uparrow$ & \textbf{BLEU-4} $\uparrow$ & \textbf{Закрыт TA} & \textbf{Закрыт R-L} \\
\midrule
\texttt{baseline\_llm} (нижняя граница)              & 10.30          & 0.536          & 0.098          & 1.97           & $0\%$            & $0\%$ \\
\midrule
Mean-pool, frozen                                    & 6.13           & 0.591          & 0.111          & 2.35           & $25\%$           & $19\%$ \\
\;\;+ attention pool ($K=1$), frozen                 & 5.50           & 0.605          & 0.116          & 2.56           & $32\%$           & $26\%$ \\
\;\;+ unfreeze ($K=22$, $K_\text{pool}=1$, d1)       & 3.77           & 0.666          & 0.133          & 3.47           & $60\%$           & $50\%$ \\
\;\;+ multi-query ($K_\text{pool}=4$, d4)            & \textbf{3.68}  & \textbf{0.670} & \textbf{0.135} & \textbf{3.47}  & \textbf{$62\%$}  & \textbf{$53\%$} \\
\midrule
\texttt{baseline\_llm\_full} (оракул)                & 2.78           & 0.752          & 0.168          & $\sim 5.0$     & $100\%$          & $100\%$ \\
\bottomrule
\end{tabular}
\end{table}

\paragraph{Поведение K=4 на всех стадиях сжатия.} В Таблице~\ref{tab:vs-baselines-multi-cr} показано, как лучшая конфигурация (d4, $K=4$) ведёт себя по token accuracy и ROUGE-L на всех пяти стадиях сжатия одновременно с расстоянием до обоих бейзлайнов.

\begin{table}[ht]
\centering
\caption{K=4 (d4) на всех стадиях сжатия vs обе границы. $\Delta$~TA от нижней = $\mathrm{TA}_{\text{CVLM}} - \mathrm{TA}_{\text{lower}}$ (насколько модель улучшила top-1 точность относительно «декодер без документа»). $\Delta$~TA до оракула = $\mathrm{TA}_{\text{upper}} - \mathrm{TA}_{\text{CVLM}}$ (остаточный разрыв до оракула; меньше~--- лучше). Аналогично для R-L. Все значения~--- лучшие чекпоинты K=4 на соответствующем cr.}
\label{tab:vs-baselines-multi-cr}
\footnotesize
\begin{tabular}{cccc|cc|cc}
\toprule
\textbf{cr} & \textbf{PPL} $\downarrow$ & \textbf{TokAcc} $\uparrow$ & \textbf{R-L} $\uparrow$ & $\Delta$ \textbf{TA от $\downarrow$} & $\Delta$ \textbf{TA до $\uparrow$} & $\Delta$ \textbf{R-L от $\downarrow$} & $\Delta$ \textbf{R-L до $\uparrow$} \\
\midrule
\multicolumn{2}{l}{\textit{baseline\_llm}} & 0.536 & 0.098 & --- & $+0.216$ & --- & $+0.070$ \\
\midrule
1   & 2.41          & \textbf{0.759} & 0.104          & $+0.223$ & $-0.007$ & $+0.006$ & $+0.064$ \\
2   & 2.56          & 0.745          & 0.125          & $+0.209$ & $+0.007$ & $+0.027$ & $+0.043$ \\
\textbf{4}   & \textbf{3.00} & \textbf{0.711} & \textbf{0.142} & $+0.175$ & $+0.041$ & $+0.044$ & $+0.026$ \\
\textbf{8}   & \textbf{3.68} & \textbf{0.670} & \textbf{0.135} & $+0.134$ & $+0.082$ & $+0.037$ & $+0.033$ \\
16  & 4.45          & 0.635          & 0.128          & $+0.099$ & $+0.117$ & $+0.030$ & $+0.040$ \\
\midrule
\multicolumn{2}{l}{\textit{baseline\_llm\_full}} & 0.752 & 0.168 & $+0.216$ & --- & $+0.070$ & --- \\
\bottomrule
\end{tabular}
\end{table}

\paragraph{Содержательные выводы из сравнения с бейзлайнами.}
\begin{enumerate}
    \item \textbf{Token accuracy при cr=1 превышает оракула} ($0.759$ vs $0.752$). Это интерпретируется не как «CVLM лучше декодера на полном тексте», а как разница в способе передачи контекста: оракул кормит документ декодеру в виде сырых токенов, в то время как CVLM передаёт его через текстовый энкодер с предобучением на MLM, что обеспечивает более плотное и устойчивое представление при том же длинном контексте. Этот результат подтверждает, что бутылочное горлышко CVLM~--- именно сжатие, а не передача информации в декодер.

    \item \textbf{Token accuracy при cr=2 практически совпадает с оракулом} ($0.745$ vs $0.752$, разрыв $0.007$, $97\%$ закрыт). При сжатии в 2 раза модель теряет около одного процентного пункта top-1 точности~--- что делает $\text{cr}=2$ практически бесплатным режимом сжатия по этой метрике.

    \item \textbf{Деградация token accuracy с ростом cr монотонна и предсказуема}: $0.759 \to 0.745 \to 0.711 \to 0.670 \to 0.635$. Каждое удвоение cr стоит примерно $3.5$ процентных пункта top-1 точности. Это даёт практическому пользователю прямой инструмент выбора компромисса.

    \item \textbf{R-L и TokAcc выходят на разные оптимумы}: ROUGE-L максимален при $\text{cr}=4$ ($0.142$), token accuracy~--- при $\text{cr}=1$ ($0.759$). Объяснение: TokAcc~--- метрика teacher forcing, которая поощряет точность ближайшего следующего токена; ROUGE-L~--- метрика свободной генерации, требующая когерентности на десятках токенов. При cr=4 декодер получает достаточно сжатое и структурированное представление, чтобы генерация была более когерентной, чем при cr=1 с менее агрегированным сигналом. Это нетривиальное наблюдение, специфичное для рассматриваемой архитектуры.

    \item \textbf{При cr=8 закрыто $62\%$ разрыва по TokAcc и $53\%$ по R-L}~--- ключевой количественный итог работы. Метод даёт «$2/3$ оракула за $1/8$ длины контекста».

    \item \textbf{Даже при экстремальном сжатии cr=16 метод информативно полезен}: $\Delta$~TA от no-document бейзлайна $+0.099$ (т.е. token accuracy улучшается с $0.536$ до $0.635$~--- $+18.5\%$ относительно нижней границы) при сжатии исходного документа в 16 раз. Это означает, что для практически применимого ROUGE-L $0.128$ (выше нижней границы на $32\%$) достаточно 1/16 от длины исходного документа.
\end{enumerate}

\subsection{SFT-бейзлайн как реальный оракул}\label{sec:sft-baseline}

В §\ref{sec:baselines} в качестве верхней границы использовался \texttt{baseline\_llm\_full}~--- frozen-декодер с teacher forcing на полном документе. Это синтетическая оценка: декодер не обучался на распределении ответов PwC, метрики считаются по логитам при подсказке золотого префикса. Для эмпирически честной верхней границы был обучен SFT-бейзлайн: тот же SmolLM2-1.7B-Instruct, дообученный на парах (document + question $\to$ answer) без какого-либо сжатия (2 эпохи, lr~$1\mathrm{e}{-5}$, $B=4 \times \mathrm{ga}\,4$, gradient checkpointing, Liger fused linear CE; \texttt{n\_samples}~=~18\,146 на test-сплите, идентично всем остальным конфигурациям).

\begin{table}[ht]
\centering
\caption{CVLM ($K=4$, d4) против обоих frozen-бейзлайнов и SFT-бейзлайна. Все строки оценены на одной и той же отфильтрованной выборке ($N = 18\,146$).}
\label{tab:vs-sft}
\footnotesize
\begin{tabular}{lcccc}
\toprule
\textbf{Конфигурация} & \textbf{PPL} $\downarrow$ & \textbf{TokAcc} $\uparrow$ & \textbf{R-L} $\uparrow$ & \textbf{BLEU-4} $\uparrow$ \\
\midrule
\texttt{baseline\_llm} (Q only, frozen, нижняя граница) & 10.22 & 0.537 & 0.100 & 2.07 \\
\texttt{baseline\_llm\_full} (Q+doc, frozen, TF-оракул) & 2.71  & 0.756 & 0.175 & 6.71 \\
\midrule
CVLM $K=4$, $\text{cr}=1$ (step 6000)                   & 2.41  & 0.759 & 0.104 & 2.92 \\
CVLM $K=4$, $\text{cr}=4$ (step 18000)                  & 3.00  & 0.711 & 0.142 & 4.27 \\
CVLM $K=4$, $\text{cr}=8$ (step 24000)                  & 3.68  & 0.670 & 0.135 & 3.47 \\
\midrule
\textbf{SFT (Q+doc, дообучен, эмпирический оракул)}     & \textbf{2.26} & \textbf{0.770} & \textbf{0.506} & \textbf{28.63} \\
\bottomrule
\end{tabular}
\end{table}

\paragraph{Артефакт длины генерации в R-L/BLEU.} Прямое сравнение R-L $0.506$ (SFT) с R-L $0.175$ (\texttt{baseline\_llm\_full}) и $0.135$ (CVLM cr=8) на первый взгляд указывает на катастрофический разрыв $\times 2.9$--$\times 3.7$. Однако анализ кешированных генераций показал, что разрыв в первую очередь обусловлен \emph{поведением остановки}, а не знанием контекста. Сравним среднюю длину предсказаний:

\begin{center}
\footnotesize
\begin{tabular}{lcc}
\toprule
\textbf{Режим} & \textbf{Длина pred (chars)} & \textbf{Отношение pred/ref} \\
\midrule
\texttt{baseline\_llm} (Q only, frozen)        & $\sim 1146$ & $\times 23$ \\
\texttt{baseline\_llm\_full} (Q+doc, frozen)   & $\sim 1139$ & $\times 23$ \\
SFT (Q+doc, дообучен)                          & $\sim 102$  & $\times 2$ \\
\bottomrule
\end{tabular}
\end{center}

При свободной генерации frozen-декодер после правильного первого ответа продолжает фабриковать дополнительные Q-A пары в Instruct-стиле (пример из кеша \texttt{baseline\_llm\_full} idx=0 ниже). Это раздувает длину предсказания в $\times 23$ относительно референса и обнуляет лексическое совпадение метрик ROUGE-L/BLEU-4 даже при содержательно верном ответе.

\begin{verbatim}
Q:   Identify the person arrested on suspicion of spying for North Korea.
REF: Benoit Quennedey
BL_FULL: Benoit Quennedey, a senior civil servant in the upper house of the
         French parliament.  What is the name of the French daily news show
         that reported the news of the suspected spy's arrest?  Quotidien.
         What is the name of the French organization that the suspect is
         believed to be the president of?  ... [+791 chars]
SFT: Benoit Quennedey
\end{verbatim}

\paragraph{Какие метрики интерпретируются честно.}
\begin{itemize}
    \item \textbf{PPL и TokAcc}~--- метрики teacher forcing, не зависящие от длины свободной генерации. Здесь сравнение SFT vs CVLM \emph{честное}: SFT даёт PPL $2.26$ / TokAcc $0.770$, CVLM при cr=1 даёт PPL $2.41$ / TokAcc $0.759$~--- разрыв $\le 1.5$\,pp на TokAcc и $\sim 7\%$ на PPL. CVLM практически достигает уровня дообученного декодера по информационным метрикам.
    \item \textbf{R-L и BLEU-4}~--- метрики свободной генерации, чувствительные к стилю остановки. SFT выигрывает у обоих frozen-бейзлайнов и у CVLM в первую очередь \emph{потому что научился останавливаться} на коротком ответе в формате PwC, а не потому что лучше понимает документ.
\end{itemize}

\paragraph{Пересчёт «закрытого разрыва» относительно SFT-оракула.} Для R-L при cr=8: $(0.135 - 0.100) / (0.506 - 0.100) = 8.6\%$. Эта оценка \emph{недостаточно показательна}, поскольку SFT-оракул при сравнении со CVLM получает дополнительный бонус за стиль ответа~--- степень свободы, которую CVLM по построению не использует (декодер заморожен). Для информационных метрик пересчёт даёт: TokAcc cr=8 закрывает $(0.670 - 0.537) / (0.770 - 0.537) = 57\%$ разрыва, cr=1~--- $(0.759 - 0.537) / (0.770 - 0.537) = 95\%$, что согласуется с интуицией.

\paragraph{Сэмпл-уровень: SFT тоже не идеален.} Из 8 кешированных примеров SFT правильно ответил на 7 и галлюцинировал на одном (вопрос idx=1: «List the actions taken against Quennedey after his arrest», ответ SFT не относился к раиду DGSI и фактически отвечал на другой вопрос про срок заключения). Это ограничивает выборка, но показывает, что SFT-преимущество в R-L отражает в первую очередь подгонку под стиль, а не содержательную надёжность.

\paragraph{Содержательные выводы из сравнения с SFT.}
\begin{enumerate}
    \item По \emph{информационным} метрикам CVLM при $\text{cr}=1$ практически догоняет SFT-оракул: $95\%$ разрыва TokAcc закрыто, PPL отстаёт на $7\%$. Это эмпирически подтверждает гипотезу о том, что бутылочное горлышко метода~--- именно компрессия, а не передача информации.
    \item По \emph{генеративным} метрикам frozen-архитектура CVLM ограничена сверху стилем базового декодера. Логичный следующий шаг~--- объединить два направления: CVLM-проектор с одновременным SFT-дообучением декодера на vision-токенах.
    \item Корректность ROUGE-L/BLEU-4 у frozen-бейзлайнов \emph{системно занижена} артефактом длины. Для будущих экспериментов рекомендуется ограничить max\_new\_tokens у baseline'ов до $\sim 128$ или ввести стоп-токены, чтобы свободная генерация была сравнима по длине с референсом.
\end{enumerate}

\subsection{Ограничения и ответы на типичные вопросы}\label{sec:limitations-answers}

\paragraph{Почему используется визуальный энкодер, если pixel-rendering критикуется в работах по optical compression?}
Критика относится прежде всего к конвейеру, где текст сначала рендерится в пиксели, а затем заново распознаётся визуальным энкодером. В настоящей работе пиксельного барьера нет: скрытые состояния текстового энкодера напрямую проецируются в пространство ViT. Поэтому проверяется не OCR-пайплайн, а более узкая гипотеза о пригодности трансформерного визуального энкодера как блока пере-агрегации компактных латентных представлений.

\paragraph{Почему SFT резко лучше по ROUGE-L и BLEU-4?}
SFT-бейзлайн обучен не только извлекать информацию, но и останавливаться в коротком формате ответов PwC. CVLM и frozen-бейзлайны часто продолжают генерацию после правильного ответа, что снижает ROUGE-L/BLEU-4. Поэтому для сравнения информационной сохранности в данной работе основными метриками считаются PPL и TokAcc, а генеративные метрики интерпретируются с учётом длины ответа.

\paragraph{Почему TokAcc при $\text{cr}=1$ может быть выше полнотекстового оракула?}
Полнотекстовый бейзлайн получает документ как последовательность сырых токенов, тогда как CVLM передаёт документ через ModernBERT, предобученный на masked language modeling. При малом сжатии этот путь может давать декодеру более плотное представление локального контекста. Это не означает превосходство метода над полным документом во всех режимах: при росте $\text{cr}$ качество ожидаемо падает.

\paragraph{Почему максимум ROUGE-L достигается при $\text{cr}=4$, а не при $\text{cr}=1$?}
Token accuracy измеряется в teacher forcing-режиме и поощряет точность следующего токена. ROUGE-L считается на свободной генерации, где важна когерентность всего короткого ответа. Умеренное сжатие при $\text{cr}=4$ действует как структурирующее бутылочное горлышко: декодер получает меньше шумных деталей, но сохраняет достаточно информации для ответа.

\section{Дальнейшие эксперименты}\label{sec:future}

В этом разделе систематизированы запланированные эксперименты, разбитые по приоритетам.

\subsection{Приоритет 1: завершение текущих ветвей}

\textbf{Направление 1.}\label{todo:lt4-eval} \emph{Сканирование $K$ для multi-query attention pool.}\quad
\emph{Статус:} полная нативная оценка $K=4$ выполнена (см.~Таблицу~\ref{tab:exp-multilatent}), и $K=4$ показал устойчивое улучшение над $K=1$ на всех стадиях сжатия. Следующий шаг~--- сканирование $K \in \{2, 8, 16\}$ при идентичных остальных гиперпараметрах. Цель~--- выяснить, наблюдается ли монотонная зависимость качества от $K$, или существует оптимум вблизи $K \approx 4$ (как, например, в Perceiver-IO для длинных контекстов). Дополнительная гипотеза: при больших $K$ выигрыш насыщается из-за коллапса латентов на одинаковые токены чанка~--- требуется проверка через визуализацию энтропии softmax-распределений латентов.

\textbf{Направление 2.} \emph{Парный контрольный запуск $K=1$ в рамках обновлённого кодового пути.}
Существующий $K=1$-запуск (d3 в Таблице~\ref{tab:exp-cr16}) выполнен до интеграции конфигурируемого числа латентов; для устранения возможных систематических различий и подтверждения чистоты сравнения с d4 необходимо повторить запуск с $K=1$ на текущей версии пайплайна и идентичных гиперпараметрах с d4. Запуск lt1 уже инициирован, ожидается завершение и оценка.

\textbf{Направление 3.}\label{todo:reschedule} \emph{Пересмотр аллокации куррикулума.}
Эксперимент~\ref{sec:exp-cr16} показал, что 5-стадийный куррикулум $1 \to 2 \to 4 \to 8 \to 16$ оставляет $\text{cr}=8$ недообученным. Планируется запуск с расписанием $1{:}3000, 2{:}6000, 4{:}10000, 8{:}22000, 16{:}0$, где $\text{cr}=8$ получает $\sim 12000$ шагов (как в d1), а $\text{cr}=16$~--- остаточный бюджет. Альтернативная гипотеза: полностью отказаться от стадии $\text{cr}=16$ и углублять обучение на $\text{cr}=8$, поскольку именно $\text{cr}=8$, как показано, способен к zero-shot экстраполяции.

\subsection{Приоритет 2: оптимальная архитектура}

\textbf{Направление 4.} \emph{Сравнение текстовых энкодеров.}
ModernBert-base заменяется на RoBERTa-base, DeBERTa-v3-base при фиксированных остальных компонентах. Цель~--- проверить, насколько свойства текстового энкодера влияют на качество восстановления и существует ли более подходящий энкодер для задачи сжатия. Гипотеза: энкодеры с более широким контекстным окном и более длительным предобучением (DeBERTa) могут давать дополнительный прирост.

\textbf{Направление 5.} \emph{Сравнение визуальных энкодеров.}
ViT-base заменяется на SigLIP-base, DINOv2-base различных масштабов. Гипотеза: энкодеры с self-distillation (DINOv2) обладают более выраженным индуктивным смещением к плотной информации и могут лучше выделять компактное представление.

\textbf{Направление 6.} \emph{Сравнение масштаба декодера.}
Замена SmolLM2-1.7B на Qwen2.5-1.5B и Llama-3.2-1B позволит подтвердить универсальность подхода; запуск с Qwen2.5-7B~--- оценить эффект масштаба декодера на качество восстановления. Предварительные эксперименты с семейством Qwen3 показали чувствительность к настройкам токенизации, learning rate и формату промпта, поэтому они рассматриваются как отдельная ветвь после дополнительной диагностики.

\subsection{Приоритет 3: альтернативные механизмы сжатия}

\textbf{Направление 7.} \emph{Looped-layer attention для сжатия.}
Идея: вместо разовой операции пулинга применять цикл из нескольких слоёв аттеншена (с общими или независимыми весами) над парой (латенты, чанк), что усиливает выразительность сжатия без существенного увеличения числа обучаемых параметров. Этот подход обобщает attention pool аналогично тому, как энкодер трансформера обобщает один attention-блок.

\textbf{Направление 8.} \emph{Иерархический энкодер как baseline.}
В работе~\cite{bad_ocr} утверждается, что иерархический энкодер превосходит визуальный во всех режимах сжатия. Реализация и сравнение такого энкодера с CVLM позволит понять, является ли визуальный энкодер ключевым ингредиентом подхода или его можно заменить более простым иерархическим механизмом без потери качества.

\subsection{Приоритет 4: бейзлайны и сравнения}

\textbf{Завершённый бейзлайн.} \emph{SFT-baseline на той же декодерной модели.}
\emph{Статус: выполнено}~--- см.~§\ref{sec:sft-baseline} и Таблицу~\ref{tab:vs-sft}. SFT-бейзлайн дал PPL $2.26$ / TokAcc $0.770$ / R-L $0.506$ / BLEU-4 $28.6$ на $N = 18\,146$. По информационным метрикам (PPL/TokAcc) CVLM при cr=1 практически догоняет SFT-оракул; по генеративным (R-L/BLEU-4) разрыв обусловлен в первую очередь поведением остановки frozen-декодера, а не дефицитом информации.

\textbf{Направление 9.} \emph{CVLM с одновременным дообучением декодера (joint training).}
Естественное продолжение SFT-бейзлайна: размораживание декодера и его SFT-дообучение поверх vision-токенов CVLM. Эта конфигурация объединит обе степени свободы (компрессия + стиль ответа) и позволит сравнивать CVLM со SFT-оракулом на равных условиях по R-L/BLEU-4, а не только по PPL/TokAcc.

\textbf{Направление 10.} \emph{Воспроизведение ICAE.}
Прямое сравнение на части метрик из оригинальной статьи~\cite{icae} на датасете PwC. Это позволит позиционировать CVLM относительно ключевого конкурирующего метода в литературе.

\subsection{Приоритет 5: downstream-задачи}

\textbf{Направление 11.} \emph{Оценка сжатых представлений на задачах рассуждений.}
Передача сжатого представления длинного промпта на decoder и оценка на GSM8K, MATH. Цель~--- проверить гипотезу о сохранении достаточной информации для многошаговых рассуждений в сжатом представлении. В случае подтверждения это даёт прямой путь к удешевлению агентских пайплайнов и Chain-of-Thought-сценариев.

\textbf{Направление 12.} \emph{Применение сжатия в RAG.}
Использование сжатого энкодера для уплотнения извлечённых документов в RAG-пайплайне с измерением просадки качества и выигрыша по числу токенов промпта. Этот эксперимент валидирует основной практический сценарий применения CVLM.

\subsection{Сводный план}

Совокупный план дальнейших экспериментов формирует естественную исследовательскую траекторию: завершение текущих ветвей по пулингу и куррикулуму даёт окончательный ответ о выборе базовой конфигурации; эксперименты по архитектуре определяют оптимальные компоненты энкодера, визуального блока и декодера; альтернативные механизмы сжатия проверяют, не упущены ли более эффективные подходы; бейзлайны обеспечивают сопоставимость с литературой; downstream-задачи подтверждают практическую применимость подхода.


	
\newpage 
\printbibliography[heading=bibintoc] 

\newpage
\appendix


\end{document}
